\markdownRendererDocumentBegin
\markdownRendererWarning{The `hybrid` option has been soft-deprecated.}{The `hybrid` option has been soft-deprecated.}{Consider using one of the following better options for mixing TeX and markdown: `contentBlocks`, `rawAttribute`, `texComments`, `texMathDollars`, `texMathSingleBackslash`, and `texMathDoubleBackslash`. For more information, see the user manual at <https://witiko.github.io/markdown/>.}{Consider using one of the following better options for mixing TeX and markdown: `contentBlocks`, `rawAttribute`, `texComments`, `texMathDollars`, `texMathSingleBackslash`, and `texMathDoubleBackslash`. For more information, see the user manual at <https://witiko.github.io/markdown/>.}\markdownRendererSectionBegin
\markdownRendererHeadingOne{越lati位ty}\markdownRendererInterblockSeparator
{}\markdownRendererSectionBegin
\markdownRendererHeadingTwo{關卡故事}\markdownRendererInterblockSeparator
{}你們被CERN製造的黑洞甩出來到外太空，並且被天狼星文明帶去比星際足球。星際足球最特別的地方是尺度非常大，因此你們要考慮相對論效應。比賽分為上下半場，上半場為戰術復刻與戰術研判，考驗你們技術動作及物理數學。下半場要針對不同的防守陣型設計不同戰術並執行，最終目標是攻破天狼星的大門。\markdownRendererInterblockSeparator
{}
\markdownRendererSectionEnd \markdownRendererSectionBegin
\markdownRendererHeadingTwo{關卡流程及場地圖}\markdownRendererInterblockSeparator
{}\markdownRendererImage{關卡流程}{../relativiskick_flowchart.png}{../relativiskick_flowchart.png}{}\markdownRendererParagraphSeparator
{}\markdownRendererImage{場地圖}{../relativiskick_layout.png}{../relativiskick_layout.png}{}\markdownRendererSoftLineBreak
{}足球場中的藍線是上半場專用，白線及邊線是下半場專用。\markdownRendererInterblockSeparator
{}
\markdownRendererSectionEnd \markdownRendererSectionBegin
\markdownRendererHeadingTwo{闖關規則}\markdownRendererInterblockSeparator
{}\markdownRendererOlBeginTight
\markdownRendererOlItemWithNumber{1}本關卡共75分鐘，前五分鐘為讀題時間，後續時間分配如關卡流程所示。\markdownRendererOlItemEnd 
\markdownRendererOlItemWithNumber{2}讀題時間結束後可打王牌或魔法牌\markdownRendererOlItemEnd 
\markdownRendererOlItemWithNumber{3}上半場建議分工：戰術復刻約2-3人，戰術研判約2-3人。\markdownRendererOlItemEnd 
\markdownRendererOlItemWithNumber{4}破壞場佈則整關0分。\markdownRendererOlItemEnd 
\markdownRendererOlEndTight \markdownRendererParagraphSeparator
{}
\markdownRendererSectionEnd \markdownRendererSectionBegin
\markdownRendererHeadingTwo{戰術復刻}\markdownRendererInterblockSeparator
{}戰術復刻流程分為兩個步驟，分別為畫跑位及跑戰術。有6個戰術，一個戰術100分，共600分。場上會有站在特定位置防守球員（都是3格高的紙箱）。\markdownRendererInterblockSeparator
{}\markdownRendererSectionBegin
\markdownRendererHeadingThree{畫跑位}\markdownRendererInterblockSeparator
{}教練區有戰術描述，閱讀戰術描述並依照跑位定義在戰術板上畫出每次接球時人的位置。畫跑位一題25分。\markdownRendererSoftLineBreak
{}畫跑位相關規定請見教練區「戰術復刻評分說明」。\markdownRendererInterblockSeparator
{}
\markdownRendererSectionEnd \markdownRendererSectionBegin
\markdownRendererHeadingThree{跑戰術}\markdownRendererInterblockSeparator
{}球員要在10秒內完成戰術描述中的戰術。跑戰術一題75分，詳細規定請見「戰術復刻評分說明」。\markdownRendererInterblockSeparator
{}
\markdownRendererSectionEnd \markdownRendererSectionBegin
\markdownRendererHeadingThree{紅牌}\markdownRendererInterblockSeparator
{}下列情形會得到紅牌一張\markdownRendererInterblockSeparator
{}\markdownRendererUlBeginTight
\markdownRendererUlItem 在戰術語言未明確書寫可以手球時，違反手球規定（詳見教練區「手球規定」）\markdownRendererUlItemEnd 
\markdownRendererUlItem 撞倒防守球員\markdownRendererSoftLineBreak
{}得到紅牌後，該題跑戰術立即結束並以零分計算。\markdownRendererUlItemEnd 
\markdownRendererUlEndTight \markdownRendererInterblockSeparator
{}
\markdownRendererSectionEnd 
\markdownRendererSectionEnd \markdownRendererSectionBegin
\markdownRendererHeadingTwo{戰術研判}\markdownRendererInterblockSeparator
{}教練區有文本題組，考點涵蓋物理與數學，共計400分。規則如下：\markdownRendererInterblockSeparator
{}\markdownRendererOlBeginTight
\markdownRendererOlItemWithNumber{1}所有文本題的答案請書寫於答案卷上\markdownRendererOlItemEnd 
\markdownRendererOlItemWithNumber{2}請將所有欲批改的答案卷於上半場結束前交給關主，否則不予計分\markdownRendererOlItemEnd 
\markdownRendererOlItemWithNumber{3}依據文本題得分，會得到等量的足球點數，可在解防守題時購買回復牌與迴轉牌（每張售價為30個足球點數）\markdownRendererOlItemEnd 
\markdownRendererOlEndTight \markdownRendererInterblockSeparator
{}其中「相對同時性與因果律」題組最重要，其最終目標為解出解防守題的「相對論規則」。若在上半場結束前有成功解出，則可進行「戰術升級」，將會降低下半場解防守題的難度。\markdownRendererInterblockSeparator
{}
\markdownRendererSectionEnd \markdownRendererSectionBegin
\markdownRendererHeadingTwo{中場休息}\markdownRendererInterblockSeparator
{}上半場結束後，會進入5分鐘的中場休息。建議可利用這段時間研究下半場解防守題的規則。此時關主也會公布以下事項：\markdownRendererInterblockSeparator
{}\markdownRendererOlBeginTight
\markdownRendererOlItemWithNumber{1}戰術復刻的得分\markdownRendererOlItemEnd 
\markdownRendererOlItemWithNumber{2}戰術研判的得分及獲得的足球點數\markdownRendererOlItemEnd 
\markdownRendererOlItemWithNumber{3}戰術升級是否成功\markdownRendererOlItemEnd 
\markdownRendererOlEndTight \markdownRendererParagraphSeparator
{}
\markdownRendererSectionEnd \markdownRendererSectionBegin
\markdownRendererHeadingTwo{解防守題}\markdownRendererInterblockSeparator
{}為回合制遊戲，一共有五題防守題，一題180分，共900分。需依照場上形勢選擇戰術，並規劃跑位來解決問題。足球場上有數名防守球員，各具不同的性質，包括瘋狗逼搶員、區域大閘、影子盯人者、路線攔截者及守門員。各題的防守球員配置不一（詳見防守題評分方式），可自行決定解題順序。另外，開始解每一題前，可告知關主欲購買回復牌及迴轉牌。\markdownRendererParagraphSeparator
{}每回合流程如下圖：\markdownRendererParagraphSeparator
{}\markdownRendererImage{解防守題流程圖}{../2nd_half_round_fig_v3.png}{../2nd_half_round_fig_v3.png}{}\markdownRendererParagraphSeparator
{}詳細規則請見「防守題評分方式」。\markdownRendererInterblockSeparator
{}\markdownRendererSectionBegin
\markdownRendererHeadingThree{紅黃牌}\markdownRendererInterblockSeparator
{}解防守題時，關主會依照犯規情節，發給紅牌或黃牌。\markdownRendererParagraphSeparator
{}下列情形會得到黃牌一張：\markdownRendererInterblockSeparator
{}\markdownRendererUlBeginTight
\markdownRendererUlItem 傳球階段時使球出界\markdownRendererUlItemEnd 
\markdownRendererUlItem 傳球階段時使防守者倒地\markdownRendererUlItemEnd 
\markdownRendererUlItem 移動階段時未依規定移動\markdownRendererUlItemEnd 
\markdownRendererUlEndTight \markdownRendererInterblockSeparator
{}得到黃牌後，場上狀態會復原至該階段開始前，並重新進行該階段。所有防守題中，每一題的黃牌數會分開累計。\markdownRendererParagraphSeparator
{}下列情形會得到紅牌一張：\markdownRendererInterblockSeparator
{}\markdownRendererUlBeginTight
\markdownRendererUlItem 違反手球規定\markdownRendererUlItemEnd 
\markdownRendererUlItem 該防守題中累計兩張黃牌\markdownRendererUlItemEnd 
\markdownRendererUlEndTight \markdownRendererInterblockSeparator
{}得到紅牌後，該防守題立刻結束並以零分計算。\markdownRendererInterblockSeparator
{}
\markdownRendererSectionEnd \markdownRendererSectionBegin
\markdownRendererHeadingThree{卡牌功能}\markdownRendererInterblockSeparator
{}\markdownRendererUlBeginTight
\markdownRendererUlItem 回復牌：回合中任何階段皆可使用，可回到上一回合的射門階段\markdownRendererUlItemEnd 
\markdownRendererUlItem 迴轉牌：得到該防守題的第二張黃牌時可使用，使用後場上狀態會復原至該階段開始前，並重新進行該階段\markdownRendererUlItemEnd 
\markdownRendererUlEndTight \markdownRendererInterblockSeparator
{}
\markdownRendererSectionEnd 
\markdownRendererSectionEnd \markdownRendererSectionBegin
\markdownRendererHeadingTwo{計分方式}\markdownRendererInterblockSeparator
{}總分 = 戰術復刻 + 戰術研判 + 解防守題 + 成就分數。滿分為2000分。\markdownRendererInterblockSeparator
{}\markdownRendererSectionBegin
\markdownRendererHeadingThree{戰術復刻}\markdownRendererInterblockSeparator
{}每個戰術畫跑位25分，跑戰術75分，共100分。6個戰術共600分。\markdownRendererInterblockSeparator
{}
\markdownRendererSectionEnd \markdownRendererSectionBegin
\markdownRendererHeadingThree{戰術研判}\markdownRendererInterblockSeparator
{}共400分，每小題配分詳見文本題題目卷。\markdownRendererInterblockSeparator
{}
\markdownRendererSectionEnd \markdownRendererSectionBegin
\markdownRendererHeadingThree{解防守題}\markdownRendererInterblockSeparator
{}若成功進球，每題防守題得分計算為\markdownRendererSoftLineBreak
{}$180 \cdot R \cdot O  \cdot C$\markdownRendererParagraphSeparator
{}回合折減係數$R$為：\markdownRendererInterblockSeparator
{}\markdownRendererUlBeginTight
\markdownRendererUlItem 戰術升級成功 = $1 - (n-4)\cdot 0.1$\markdownRendererUlItemEnd 
\markdownRendererUlItem 戰術升級失敗 = $1 - (n-4)\cdot 0.2$\markdownRendererUlItemEnd 
\markdownRendererUlItem 大於1以1計，小於0以0計\markdownRendererUlItemEnd 
\markdownRendererUlEndTight \markdownRendererInterblockSeparator
{}越位係數$O$為：\markdownRendererInterblockSeparator
{}\markdownRendererUlBeginTight
\markdownRendererUlItem 所有觀察者不越位 = 1\markdownRendererUlItemEnd 
\markdownRendererUlItem 靜止觀察者不越位 = 0.6\markdownRendererUlItemEnd 
\markdownRendererUlItem 越位 = 0.2\markdownRendererUlItemEnd 
\markdownRendererUlEndTight \markdownRendererInterblockSeparator
{}因果律係數$C$為：\markdownRendererInterblockSeparator
{}\markdownRendererUlBeginTight
\markdownRendererUlItem 所有傳球皆符合因果律 = 1\markdownRendererUlItemEnd 
\markdownRendererUlItem 有至少一次傳球違反因果律 = 0.5\markdownRendererUlItemEnd 
\markdownRendererUlEndTight \markdownRendererInterblockSeparator
{}每題最低零分。若獲得紅牌或未進球，該防守題零分。\markdownRendererInterblockSeparator
{}
\markdownRendererSectionEnd \markdownRendererSectionBegin
\markdownRendererHeadingThree{成就分數}\markdownRendererInterblockSeparator
{}|成就名稱|內容|分數|\markdownRendererSoftLineBreak
{}|-|-|-|\markdownRendererSoftLineBreak
{}|終場清道夫|協助關主場復|50分|\markdownRendererSoftLineBreak
{}|全隊到齊|全隊以足球先發陣容的形式合影，並手持紀念旗（可向關主索取範例照片）|50分|
\markdownRendererSectionEnd 
\markdownRendererSectionEnd 
\markdownRendererSectionEnd \markdownRendererDocumentEnd